\section{deal gaz. Gaz MKN asosiy tenglamasi. Dalton qonuni}


\begin{enumerate}

\item Gazni ideal deb hisoblash uchun nimani hisobga olmaslik kerak?
\begin{enumerate}[label=\Alph*)]
\item molekulalarning to‘qnashuvini
\item molekulalarning to‘qnashuvidagi o‘zaro ta’sirini
\item molekulalarning harakatini
\item molekulalarning masofadan ta’sirlashishini
\end{enumerate}

\item $\overline{\vartheta_x^2} = \frac{1}{3} \overline{\vartheta^2}$ ifoda quyida keltirilgan mulohazalarning qaysi biriga asoslanib yozilgan?
\begin{enumerate}[label=\Alph*)]
\item molekulalar soni kam
\item molekulalar o‘zaro elastik to‘qnashadi
\item molekulalar shar shaklida
\item molekulalar tartibsiz harakat qiladi
\end{enumerate}

\item Ideal gazning bosimi 166 kPa, molekulalar konsentratsiyasi $1,14 \cdot 10^{26} \, m^{-3}$ va bitta molekula massasi $2,8 \cdot 10^{-26} \, kg$. Molekulalarning o‘rtacha kvadrat tezligini toping (m/s).
\begin{enumerate}[label=\Alph*)]
\item 845
\item 220
\item 395
\item 350
\end{enumerate}

\item Ideal gazning bosimi 266 kPa, molekulalar konsentratsiyasi $1,14 \cdot 10^{26} \, m^{-3}$ va bitta molekula massasi $2,8 \cdot 10^{-26} \, kg$. Molekulalarning o‘rtacha kvadrat tezligini toping (m/s).
\begin{enumerate}[label=\Alph*)]
\item 845
\item 220
\item 400
\item 500
\end{enumerate}

\item Bitta molekulasining o‘rtacha kinetik energiyasi $7,5 \cdot 10^{-21} \, J$ va molekulalar konsentratsiyasi $1,14 \cdot 10^{24} \, m^{-3}$ bo‘lgan gazning bosimini toping (kPa).
\begin{enumerate}[label=\Alph*)]
\item 84,5
\item 22
\item 40
\item 5,7
\end{enumerate}

\item Bitta molekulasining o‘rtacha kinetik energiyasi $8,5 \cdot 10^{-21} \, J$ va molekulalar konsentratsiyasi $1,14 \cdot 10^{24} \, m^{-3}$ bo‘lgan gazning bosimini toping (kPa).
\begin{enumerate}[label=\Alph*)]
\item 84,55
\item 22,11
\item 40
\item 6,46
\end{enumerate}

\item Molekulalarning o‘rtacha kvadratik tezligi 1 km/s va molekulalarning massasi $3 \cdot 10^{-27} \, kg$ bo‘lganda $1 \, mm^3$ da 100 milliard bo‘lgan konsentratsiyadagi ideal gaz idish devorlariga qanday bosim (Pa) beradi?
\begin{enumerate}[label=\Alph*)]
\item 0,1
\item 0,8
\item 3
\item 0,5
\end{enumerate}

\item Agar 1-gaz molekulalarining massasi 2-gaz molekulalarining massasidan 4 marta katta bo‘lsa, 1-gaz molekulalarining konsentratsiyasi esa 2-gaznikidan 2 marta kichik bo‘lsa, gaz zichliklarining nisbati $\rho_1 / \rho_2$ qanday bo‘ladi?
\begin{enumerate}[label=\Alph*)]
\item 2
\item 4
\item 1/8
\item 1/4
\end{enumerate}

\item Gaz molekulasining o‘rtacha kvadratik tezligi 500 m/s va gazning bosimi 121 kPa bo‘lsa, gaz zichligini toping ($kg/m^3$).
\begin{enumerate}[label=\Alph*)]
\item 4,55
\item 2,11
\item 4
\item 1,45
\end{enumerate}

\item Vektor kattaliklarni ko‘rsating: 1) modda miqdori; 2) molekulalar tezligi; 3) molekulalar konsentratsiyasi; 4) molekulalar bosimi.
\begin{enumerate}[label=\Alph*)]
\item 2, 4
\item 1, 3
\item 2
\item 3, 4
\end{enumerate}

\item Agar kislorodning bosimi 0,2 MPa molekulalarining o‘rtacha kvadratik tezligi 700 m/s ga teng bo‘lsa, uning molekulalari konsentratsiyasini toping ($m^{-3}$).
\begin{enumerate}[label=\Alph*)]
\item $2,3 \cdot 10^{24}$
\item $2,3 \cdot 10^{25}$
\item $2,3 \cdot 10^{27}$
\item $2,8 \cdot 10^{25}$
\end{enumerate}

\item Gaz 6 kg massaga ega. U 200 kPa bosimda $5 \, m^3$ hajmni egallasa, shu gaz molekulalari harakatining o‘rtacha kvadratik tezligi qanday bo‘ladi (m/s)?
\begin{enumerate}[label=\Alph*)]
\item 707
\item 70
\item 210
\item 240
\end{enumerate}

\item Gaz molekulasining o‘rtacha kvadratik tezligi 370 m/s va gazning bosimi 66,2 kPa bo‘lsa, gaz zichligini toping ($kg/m^3$).
\begin{enumerate}[label=\Alph*)]
\item 4,55
\item 2,11
\item 4
\item 1,45
\end{enumerate}

\item Quyidagi qaysi munosabat molekulyar kinetik nazariyaning asosiy tenglamasidir:
1) $p = \frac{1}{3} n k T$; 2) $pV = mRT/\mu$; 3) $p = (1/3) n m_0 \overline{v^2}$; 4) $E_k = 3kT/2$?
\begin{enumerate}[label=\Alph*)]
\item 1
\item 2
\item 3
\item 4
\end{enumerate}

\item Bir atomli gazning hajmi ikki marta orttirilib, molekulalar o‘rtacha kinetik energiyasi ham ikki marta orttirilsa, gazning bosimi qanday o‘zgaradi?
\begin{enumerate}[label=\Alph*)]
\item 4 marta ortadi
\item 4 marta kamayadi
\item 2 marta ortadi
\item o‘zgarmaydi
\end{enumerate}

\item Agar gazning hajmi 2 marta ortsa va gaz molekulalarining o‘rtacha kvadratik tezligi shuncha marta kamaysa, uning bosimi qanday o‘zgaradi?
\begin{enumerate}[label=\Alph*)]
\item 8 marta kamayadi
\item 4 marta ortadi
\item 4 marta kamayadi
\item 8 marta ortadi
\end{enumerate}

\item Gazning bosimi 6 marta orttirilib, hajmi 1,5 marta kamaytirilsa, gaz molekulalarining o‘rtacha kvadratik tezligi qanday o‘zgaradi?
\begin{enumerate}[label=\Alph*)]
\item 2 marta ortadi
\item 2 marta kamayadi
\item 4 marta ortadi
\item 4 marta kamayadi
\end{enumerate}

\item Bir atomli ideal gaz bosimi 250 kPa molekulalarining o‘rtacha kinetik energiyasi $7,5 \cdot 10^{-21} \, J$ ga teng. Gaz molekulalarining konsentratsiyasi nimaga teng?
\begin{enumerate}[label=\Alph*)]
\item $5 \cdot 10^{25}$
\item $5 \cdot 10^{24}$
\item $5 \cdot 10^{23}$
\item $5 \cdot 10^{22}$
\end{enumerate}

\item Berk idishdagi ideal gazning molekulalari o‘rtacha kvadratik tezligi 20\% ga ko‘paytirilsa, gazning bosimi necha foizga o‘zgaradi?
\begin{enumerate}[label=\Alph*)]
\item 22\% ga kamayadi
\item 22\% ga ko‘payadi
\item 44\% ga ko‘payadi
\item 44\% ga kamayadi
\end{enumerate}

\item Tezligi 1200 m/s bo‘lgan kislorod molekulasining idish devoriga $60^\circ$ burchak ostida mutlaq elastik urulishi natijasida idish devori olgan kuch impulsini toping ($N \cdot s$). $N_A = 6 \cdot 10^{23} \, mol^{-1}$.
\begin{enumerate}[label=\Alph*)]
\item $3,2 \cdot 10^{-23}$
\item 0
\item $3,2 \cdot 10^{-25}$
\item $6,4 \cdot 10^{-23}$
\end{enumerate}

\end{enumerate}



\begin{table}[h!]
\centering
\begin{tabular}{|c|c|c|c|c|c|c|c|c|c|}
\hline
№ & J. & № & J. & № & J. & № & J. & № & J. \\ \hline
1 & D & 5 & D & 9 & D & 13 & D & 17 & A \\ \hline
2 & D & 6 & D & 10 & C & 14 & C & 18 & A \\ \hline
3 & C & 7 & A & 11 & B & 15 & D & 19 & C \\ \hline
4 & D & 8 & A & 12 & A & 16 & A & 20 & D \\ \hline
\end{tabular}
\end{table}